\documentclass[11pt]{article}
% ===========================================================================
%  Shared preamble for the AP Calculus AB exam-preparation worksheets.
%  Input right after \documentclass. Files beginning with "_" are not compiled
%  on their own. Navy and gold, to match the course's lesson pages.
% ===========================================================================
\usepackage[letterpaper,top=2.3cm,bottom=2.2cm,left=1.7cm,right=1.7cm,headheight=28pt,footskip=30pt]{geometry}
\usepackage{amsmath,amssymb}
\usepackage[T1]{fontenc}
\usepackage[utf8]{inputenc}
\usepackage{lmodern}
\usepackage{xcolor}
\usepackage[most]{tcolorbox}
\usepackage{enumitem}
\usepackage{fancyhdr}
\usepackage{array}
\usepackage{tikz}
\usepackage{pgfplots}
\pgfplotsset{compat=1.18}

\definecolor{navy}{HTML}{161E54}
\definecolor{ink}{HTML}{27357E}
\definecolor{gold}{HTML}{B98A00}
\definecolor{goldbg}{HTML}{FFFDF3}
\definecolor{inkbg}{HTML}{F6F7FD}
\definecolor{rule}{HTML}{C9CFF0}
\definecolor{okgreen}{HTML}{1B7A44}
\definecolor{okbg}{HTML}{F1F8F3}

\pagestyle{fancy}
\fancyhf{}
\fancyhead[L]{\footnotesize Name: \underline{\hspace{4cm}}}
\fancyhead[C]{\textbf{\color{navy}Integration Academy}\\[-2pt]{\scriptsize Calculus AB \textperiodcentered{} AP\textsuperscript{\textregistered} Exam Preparation}}
\fancyfoot[L]{\tiny\color{black!55}\parbox[t]{0.78\textwidth}{AP\textsuperscript{\textregistered} is a registered trademark of the College Board, which was not involved in the production of, and does not endorse, this material. \textcopyright{} Integration Academy.}}
\fancyfoot[R]{\scriptsize\thepage}
\renewcommand{\headrulewidth}{1.2pt}
\renewcommand{\headrule}{\hbox to\headwidth{\color{gold}\leaders\hrule height \headrulewidth\hfill}}
\renewcommand{\footrulewidth}{0pt}
\setlength{\parindent}{0pt}
\setlength{\parskip}{4pt}

% Title block: \wstitle{code}{title}{unit line}
\newcommand{\wstitle}[3]{\begin{center}{\small\bfseries\color{gold}WORKSHEET #1}\\[3pt]{\LARGE\bfseries\color{navy}#2}\\[3pt]{\small\color{black!60}#3}\end{center}\vspace{2pt}}
% Section heading
\newcommand{\sech}[1]{\par\vspace{10pt}{\large\bfseries\color{navy}#1}\par\nobreak\vspace{2pt}{\color{rule}\hrule height0.8pt}\vspace{7pt}}
% A level tag printed beside a group of questions
\newcommand{\level}[1]{\par\vspace{6pt}{\footnotesize\bfseries\color{gold}\MakeUppercase{#1}}\par\nobreak\vspace{2pt}}

\newtcolorbox{ideabox}{colframe=ink,colback=inkbg,coltitle=white,colbacktitle=ink,fonttitle=\bfseries\small,title={Key ideas},arc=2pt,boxrule=0.8pt,left=8pt,right=8pt}
\newtcolorbox{example}[1]{colframe=rule,colback=white,coltitle=navy,colbacktitle=inkbg,fonttitle=\bfseries\small,title={#1},arc=2pt,boxrule=0.8pt,borderline west={3pt}{0pt}{ink},left=10pt,right=8pt}
\newtcolorbox{frbox}[1]{colframe=gold,colback=goldbg,coltitle=white,colbacktitle=gold,fonttitle=\bfseries\small,title={#1},arc=2pt,boxrule=0.8pt,breakable,left=8pt,right=8pt}
\newtcolorbox{solbox}{colframe=okgreen,colback=okbg,coltitle=white,colbacktitle=okgreen,fonttitle=\bfseries\small,title={Worked solutions},arc=2pt,boxrule=0.8pt,breakable,left=8pt,right=8pt}
\newtcolorbox{scorebox}[1]{colframe=navy,colback=white,coltitle=white,colbacktitle=navy,fonttitle=\bfseries\small,title={#1},arc=2pt,boxrule=0.8pt,breakable,left=8pt,right=8pt}

\newcommand{\soln}{\par\smallskip\textbf{\color{ink}Solution.}\ }
\newcommand{\result}[1]{\par\smallskip\textbf{#1}}
% blank working space of a given height
\newcommand{\work}[1]{\par\vspace{#1}}
% a point value in the margin of a scoring guide
\newcommand{\pts}[1]{\hfill{\footnotesize\bfseries\color{gold}[#1]}}

% Coordinate grid for graph questions: \begin{apgraph}{xmin}{xmax}{ymin}{ymax} ... \end{apgraph}
\newenvironment{apgraph}[4]{\begin{center}\begin{tikzpicture}\begin{axis}[
  axis lines=middle, xmin=#1, xmax=#2, ymin=#3, ymax=#4,
  grid=both, grid style={black!12}, xtick distance=1, ytick distance=1,
  ticklabel style={font=\scriptsize}, xlabel={$x$}, ylabel={$y$},
  width=10cm, height=7.2cm, samples=80, clip=false,
  every axis plot/.append style={line width=1.1pt, color=ink}]}{\end{axis}\end{tikzpicture}\end{center}}
\newcommand{\closeddot}[2]{\addplot[only marks, mark=*, mark size=2.2pt, color=ink] coordinates {(#1,#2)};}
\newcommand{\opendot}[2]{\addplot[only marks, mark=*, mark size=2.2pt, color=ink, mark options={fill=white}] coordinates {(#1,#2)};}

% Multiple-choice options: four across, or stacked when the options are long.
\newcommand{\opts}[4]{\par\vspace{3pt}\begin{tabular}{@{}p{0.235\linewidth}p{0.235\linewidth}p{0.235\linewidth}p{0.235\linewidth}@{}}(A)\ #1 & (B)\ #2 & (C)\ #3 & (D)\ #4\end{tabular}\par\vspace{9pt}}
\newcommand{\optsv}[4]{\par\vspace{3pt}\begin{tabular}{@{}l@{\ \ }p{0.9\linewidth}@{}}(A) & #1\\(B) & #2\\(C) & #3\\(D) & #4\end{tabular}\par\vspace{9pt}}

\fancyhead[R]{\footnotesize\color{black!60}Worksheet 7.3}
\begin{document}

\wstitle{7.3}{Drawing Slope Fields}{Unit 7: Differential Equations}

\begin{ideabox}
\begin{itemize}[leftmargin=*,itemsep=2pt,topsep=2pt]
  \item A \textbf{slope field} shows, at each point $(x,y)$, a short segment whose slope is the value of $\dfrac{dy}{dx}$ there.
  \item To draw one, substitute each point into the differential equation.
  \item Segments are horizontal wherever $\dfrac{dy}{dx}=0$.
  \item If $\dfrac{dy}{dx}$ depends only on $x$, every column of the field is identical; if only on $y$, every row is.
\end{itemize}
\end{ideabox}

\sech{Worked examples}
\begin{example}{Example 1 \textperiodcentered{} Computing slopes}
For $\dfrac{dy}{dx}=y-2x$, find the slopes at $(1,1)$, $(0,2)$ and $(2,4)$.\soln Substitute each point.\result{$-1$, $2$ and $0$.}
\end{example}

\begin{example}{Example 2 \textperiodcentered{} Horizontal segments}
Where are the segments of the slope field for $\dfrac{dy}{dx}=y-2x$ horizontal?\soln Set $y-2x=0$.\result{Along the line $y=2x$.}
\end{example}

\begin{example}{Example 3 \textperiodcentered{} Recognising a pattern}
A slope field has the same segments all the way up each vertical line. What does this say about the differential equation?\soln The slope does not change when $y$ changes.\result{$\dfrac{dy}{dx}$ depends on $x$ only.}
\end{example}

\newpage
\sech{Practice}

\level{Core}
\begin{enumerate}[leftmargin=*,itemsep=2pt]
  \item For $\dfrac{dy}{dx}=x^2y$, find the slope at $(2,3)$.\work{2.3cm}
  \item For $\dfrac{dy}{dx}=x^2y$, find the slope at $(-1,-2)$.\work{2.3cm}
  \item For $\dfrac{dy}{dx}=x^2y$, where are the segments horizontal?\work{2.3cm}
  \item For $\dfrac{dy}{dx}=x^2y$, what is the sign of the slopes in the third quadrant?\work{2.3cm}
\end{enumerate}

\level{Exam level}
\begin{enumerate}[leftmargin=*,itemsep=2pt,start=5]
  \item For $\dfrac{dy}{dx}=x^2+y^2$, explain why no segment has negative slope, and find where the slope is zero.\work{2.8cm}
  \item Describe the slope field of $\dfrac{dy}{dx}=(y-1)(y+2)$: where are the segments horizontal, and what is special about each row?\work{2.8cm}
  \item Describe the slope field of $\dfrac{dy}{dx}=x-2$, and the shape of its solution curves.\work{2.8cm}
  \item A slope field has horizontal segments along the curve $y=x^3$, positive slopes above that curve and negative slopes below it. Suggest a differential equation.\work{2.8cm}
\end{enumerate}

\begin{frbox}{Question 9 \textperiodcentered{} Free-response practice}
Consider the differential equation $\dfrac{dy}{dx}=y(2-x)$.
\begin{enumerate}[label=(\alph*),leftmargin=*,itemsep=2pt]
  \item Find the slope at each of the points $(0,1)$, $(2,3)$, $(3,2)$ and $(1,-1)$.\work{2.6cm}
  \item Describe all points at which the segments of the slope field are horizontal.\work{2.6cm}
  \item The solution through $(2,3)$ has a horizontal tangent there. Is this a relative maximum or a relative minimum? Justify your answer.\work{2.6cm}
\end{enumerate}
\end{frbox}

\newpage
\begin{solbox}
\begin{enumerate}[leftmargin=*,itemsep=4pt]
  \item $12$.
  \item $-2$.
  \item Along both axes: $x=0$ or $y=0$.
  \item Negative ($x^2>0$, $y<0$).
  \item $x^2+y^2\ge0$ always; it is zero only at the origin.
  \item Horizontal along $y=1$ and $y=-2$. The slope depends only on $y$, so each row has identical segments. Slopes are negative between the two lines and positive outside them.
  \item The slope depends only on $x$: columns are identical, negative for $x<2$ and positive for $x>2$. Solutions are parabolas $y=\dfrac{x^2}{2}-2x+C$ with their lowest points at $x=2$.
  \item $\dfrac{dy}{dx}=y-x^3$.
  \item \begin{enumerate}[label=(\alph*),leftmargin=*,itemsep=3pt]
    \item $2$, $0$, $-2$ and $-1$.
    \item The line $y=0$ and the line $x=2$.
    \item $\dfrac{d^2y}{dx^2}=\dfrac{dy}{dx}(2-x)-y$, which equals $-3$ at $(2,3)$. A zero first derivative and a negative second derivative: a relative maximum.
  \end{enumerate}
\end{enumerate}
\end{solbox}

\end{document}
